\chapter{Linked Lists II: Circular Linked Lists}\label{ch:circular-list}

\chapterguide%
  {Understand the node/head model and singly linked-list traversal from \chapref{ch:linked-lists}.}
  {explain the circular link, compare singly/doubly/circular lists, trace circular traversal, and implement beginning/end insertion and deletion plus deletion by value.}

A circular linked list changes one crucial link: the last node points back to the first node instead of pointing to \texttt{NULL}. The result is a sequence that can continue in a loop.

\section{The circular structure}

\begin{adsfig}[A circular linked list]{A circular linked list drawn as the ring it actually is. The single change
from \figref{fig:sll-anatomy} is the last node's \texttt{next} field: instead of
\texttt{NULL} it stores the address of the head, so there is no end to stop
at.}{fig:cll-ring}
\begin{tikzpicture}[font=\small]
% ring of four nodes
\foreach \v/\ang [count=\i from 0] in {10/90, 20/0, 30/270, 40/180}{
  \node[adsnode,minimum size=9.5mm] (r\i) at (\ang:1.75) {\v};
}
\draw[adsarrow] (r0) to[bend left=22] (r1);
\draw[adsarrow] (r1) to[bend left=22] (r2);
\draw[adsarrow] (r2) to[bend left=22] (r3);
\draw[-{Stealth[length=2.4mm]},line width=1pt,draw=ADSOrange] (r3) to[bend left=22] (r0);
\node[adstag,anchor=south] at (90:2.45) {HEAD};
\node[adslabel,anchor=east,color=ADSOrange,align=right] at (180:2.5)
  {LAST\\\texttt{last->next}\\$=$ head};
\node[adslabel,align=center] at (0,0) {no \texttt{NULL}\\anywhere};

% linear redrawing for comparison
\begin{scope}[xshift=6.9cm]
\foreach \v [count=\i from 0] in {10,20,30,40}{
  \node[adscell,minimum width=0.95cm,minimum height=0.8cm] (s\i) at (1.05*\i,0.55) {\v};
}
\foreach \i in {0,1,2}{
  \pgfmathtruncatemacro{\j}{\i+1}
  \draw[adsptr] (s\i)--(s\j);
}
\draw[-{Stealth[length=2.2mm]},line width=1pt,draw=ADSOrange]
  (s3.south) -- ++(0,-0.62) -| (s0.south);
\node[adslabel,anchor=north,align=center] at (1.57,-0.85)
  {the same list drawn flat:\\the wrap link is what makes it circular};
\end{scope}
\end{tikzpicture}
\end{adsfig}

The Week IV slides emphasize four consequences:

\begin{itemize}
  \item the last node points back to the head;
  \item there is no \texttt{NULL} at the end;
  \item traversal can continue indefinitely around the loop;
  \item therefore traversal needs an explicit stopping condition.
\end{itemize}

\section{Traversal: stop when you return to the head}

A normal singly linked list often uses \texttt{current != nullptr} as the loop condition. That condition does not work for a non-empty circular list because the pointer never becomes null.

\begin{lstlisting}
void display(Node* head) {
    if (head == nullptr) return;

    Node* current = head;
    do {
        cout << current->data << ' ';
        current = current->next;
    } while (current != head);

    cout << '\n';
}
\end{lstlisting}

\begin{keyidea}
For a non-empty circular list, the natural traversal invariant is ``keep moving until the pointer returns to the starting node.''
\end{keyidea}

That invariant explains the unusual loop shape. The test \texttt{current != head} is
\emph{already false} on the first iteration, because \texttt{current} starts at the head.
A plain \texttt{while} loop would therefore print nothing at all. The \texttt{do}\,--\,\texttt{while}
form checks the condition only after the body has run once, which is exactly what a ring
needs (\figref{fig:cll-traverse}).

\begin{adsfig}[Why circular traversal uses do--while]{Why circular traversal uses \texttt{do}\,--\,\texttt{while}. The pointer
starts on the head and returns to it, so the stopping condition is true at both the
beginning and the end of the walk; only the second occurrence should end the
loop.}{fig:cll-traverse}
\begin{tikzpicture}[font=\small]
\foreach \v/\ang [count=\i from 0] in {10/90, 20/0, 30/270, 40/180}{
  \node[adsnode,minimum size=9mm] (t\i) at (\ang:1.6) {\v};
}
\draw[adsarrow] (t0) to[bend left=22] (t1);
\draw[adsarrow] (t1) to[bend left=22] (t2);
\draw[adsarrow] (t2) to[bend left=22] (t3);
\draw[adsarrow] (t3) to[bend left=22] (t0);
\node[adstag,anchor=south] at (90:2.3) {start \& stop here};
\node[adslabel,align=center,color=ADSMuted] at (0,0) {visit order\\10, 20, 30, 40};

\begin{scope}[xshift=5.6cm,yshift=0.35cm]
\node[adsstep,fill=ADSPaleRed,draw=ADSRed,text width=5.3cm,align=left,anchor=north west]
  (bad) at (0,0.9)
  {\textcolor{ADSRed}{\bfseries Wrong}\\[1pt]
   \texttt{while (current != head) \{...\}}\\
   condition false immediately $\rightarrow$ prints nothing};
\node[adsstep,fill=ADSPaleTeal,draw=ADSGreen,text width=5.3cm,align=left,anchor=north west]
  (good) at (0,-0.75)
  {\textcolor{ADSGreen}{\bfseries Right}\\[1pt]
   \texttt{do \{...\} while (current != head);}\\
   body runs first, then the test ends the lap};
\end{scope}
\end{tikzpicture}
\end{adsfig}

\section{Singly vs doubly vs circular}

\begin{mltable}
\begin{tabularx}{\linewidth}{@{}>{\bfseries\leavevmode\color{ADSNavy}}p{0.19\linewidth}X X X@{}}
\toprule
\rowcolor{ADSPaleBlue!92}
Feature & Singly & Doubly & Circular \\
\midrule
Last node points to & \texttt{NULL} & \texttt{NULL} & head \\
Direction & Forward & Forward and backward & Usually forward circular in this course \\
Pointers per node & \texttt{next} & \texttt{prev}, \texttt{next} & \texttt{next} \\
Traversal stop & \texttt{NULL} & \texttt{NULL} & back at head \\
Main use in slides & Simple dynamic list & Two-way navigation & Repeated/cyclic processing \\
\bottomrule
\end{tabularx}
\end{mltable}

\section{Insertion at the beginning}

Assume \texttt{last} points to the final node. Then \texttt{last->next} is the head. This representation makes both the tail and head easy to identify.

\begin{mlsteps}
  \item Allocate the new node.
  \item If the list is empty, point the node to itself and make it \texttt{last}.
  \item Otherwise, point the new node to the current head.
  \item Set \texttt{last->next} to the new node, making it the new head.
\end{mlsteps}

\begin{lstlisting}
void insertBeginning(Node*& last, int value) {
    Node* n = new Node{value, nullptr};

    if (last == nullptr) {
        n->next = n;
        last = n;
        return;
    }

    n->next = last->next;
    last->next = n;
}
\end{lstlisting}

\section{Insertion at the end}

With the same \texttt{last} representation, insert at the beginning of the circle and then move \texttt{last} to the new node.

\begin{lstlisting}
void insertEnd(Node*& last, int value) {
    Node* n = new Node{value, nullptr};

    if (last == nullptr) {
        n->next = n;
        last = n;
        return;
    }

    n->next = last->next;  // new node points to head
    last->next = n;        // old last points to new node
    last = n;              // new node becomes last
}
\end{lstlisting}

\section{Deleting from the beginning}

There are three cases: empty list, one-node list, and a list with at least two nodes.

\begin{lstlisting}
bool deleteBeginning(Node*& last) {
    if (last == nullptr) return false;

    Node* head = last->next;
    if (head == last) {
        delete head;
        last = nullptr;
        return true;
    }

    last->next = head->next;
    delete head;
    return true;
}
\end{lstlisting}

\section{Deleting from the end}

A singly circular list has no direct previous pointer, so deletion of the last node requires locating the node whose \texttt{next} is \texttt{last}.

\begin{lstlisting}
bool deleteEnd(Node*& last) {
    if (last == nullptr) return false;

    Node* head = last->next;
    if (head == last) {
        delete last;
        last = nullptr;
        return true;
    }

    Node* current = head;
    while (current->next != last) {
        current = current->next;
    }

    current->next = head;
    delete last;
    last = current;
    return true;
}
\end{lstlisting}

\section{Deletion by value}

The Week IV lab also requires deletion by value. Traversal must both detect the target and recognize when the full circle has been checked.

\begin{lstlisting}
bool deleteValue(Node*& last, int value) {
    if (last == nullptr) return false;

    Node* prev = last;
    Node* current = last->next;  // head

    do {
        if (current->data == value) {
            if (current == prev) {      // one-node list
                delete current;
                last = nullptr;
            } else {
                prev->next = current->next;
                if (current == last) last = prev;
                delete current;
            }
            return true;
        }
        prev = current;
        current = current->next;
    } while (current != last->next);

    return false;
}
\end{lstlisting}

\begin{pitfall}
The easiest circular-list bug is an infinite traversal. Every loop over the list must make progress and must have a condition tied to returning to the starting node.
\end{pitfall}

\begin{quickcheck}
Why does \texttt{while (current != nullptr)} fail as a normal stopping condition in a circular list? What special pointer updates are needed when the list contains only one node?
\end{quickcheck}

\begin{chapterrecap}
\begin{itemize}
  \item A circular linked list connects the last node back to the head.
  \item The list does not terminate with \texttt{NULL}.
  \item Traversal must stop when the pointer returns to the start.
  \item Insertions and deletions require careful maintenance of the closing link.
  \item Circular linked lists are useful for repeated or cyclic processing.
\end{itemize}
\end{chapterrecap}

\labsection{4}{Circular linked-list operations}{sec:lab04}

\begin{sourcealignment}{Official Lab Manual --- Lab IV}
\textbf{Objective.} Study linked-list types more deeply and implement algorithms for circular linked lists.

\textbf{Outcomes.} Understand circular linked lists and implement one from scratch.

\textbf{Problem directions.} Implement the different insertion/deletion operations possible in a circular linked list; compare singly, doubly, and circular lists.
\end{sourcealignment}

\begin{labproject}{Maintain a circular list without breaking the loop}
\textbf{Coverage.} Chapter 4.\\
\textbf{Goal.} Maintain a circular list through every insertion and deletion without breaking the ring.

Implement insertion at beginning/end, deletion from beginning/end, deletion by value, and display.

\begin{lstlisting}
void displayCircular(Node* last) {
    if (last == nullptr) return;
    Node* head = last->next;
    Node* p = head;
    do {
        cout << p->data << ' ';
        p = p->next;
    } while (p != head);
}
\end{lstlisting}

\begin{tryit}
Use the same input sequence in singly, doubly, and circular implementations. Compare traversal stop conditions and the pointer fields required per node.
\end{tryit}
\end{labproject}

\begin{verification}
\begin{itemize}
  \item Empty, one-node, and multi-node cases are tested.
  \item After every mutation, the last node still points to the head.
  \item Traversal stops when it returns to the starting node rather than waiting for \texttt{nullptr}.
  \item Deletion-by-value handles the head, a middle node, and the last node.
\end{itemize}
\end{verification}

\begin{labdeliverable}
Submit the required operations, test output, and a three-column comparison of singly, doubly, and circular linked lists.
\end{labdeliverable}

